%=============================================================================
% P7 - Ingenius. Manuscrito maquetado en la plantilla oficial de la revista.
%
% Compilar:  pdflatex articulo_p7 && bibtex articulo_p7 && pdflatex articulo_p7 x2
%
% NOTAS DE MAQUETACION
% - El articulo esta escrito en INGLES. Las normas piden el titulo en el idioma
%   del articulo en la primera linea y la traduccion en la segunda. En la clase,
%   \titulo se imprime primero y \tituloa despues, de modo que aqui \titulo
%   lleva el INGLES y \tituloa el ESPANOL. Lo mismo con el resumen: Abstract a
%   la izquierda, Resumen a la derecha.
% - Las palabras clave van en ORDEN ALFABETICO, como pide la plantilla.
% - Las tablas llevan el titulo ARRIBA (\captionof antes del tabular).
% - ANONIMIZADO para revision doble ciego: sin nombres, adscripciones ni ORCID.
%   Rellenar \autores y \adscripcionautor solo en la version no anonima.
%=============================================================================
\documentclass[12pt,twoside,titlepage]{ingenius}
\usepackage[figuresright]{rotating}
\usepackage{pgfplots}
\pgfplotsset{compat=1.7}
\usepackage{adjustbox}
\usepackage{units}
\usepackage{dblfloatfix}
\usepackage{amssymb}
\usepackage{caption}
\usepackage{subcaption}
\usepackage{bm}
\usepackage{textcomp}
\usepackage{pifont}
\usepackage{cite}
\usepackage{xcolor,colortbl}
\usepackage{amsmath}
\usepackage{upgreek}
\usepackage{graphicx}
\usepackage{tikz}
\usetikzlibrary{shapes,arrows}
\usepackage{multirow}
\usepackage{float}
\usepackage{soul}
\usepackage{rotating}
\usepackage{hyperref}
\usepackage{booktabs}
\graphicspath{{figures/}}

%========================================================
% Titulo: ingles primero (idioma del articulo), espanol despues.
% Titulos de hasta 80 caracteres, como fija la plantilla oficial de la carta de
% presentacion de Ingenius (71 y 77). Deben coincidir con los de
% paper/llenar_carta_ingenius.py.
\titulo{Signal representation and load transfer in bearing fault
classification}
\tituloa{Representación de señal y transferencia entre cargas en fallas de
rodamientos}
%========================================================
% VERSION ANONIMIZADA PARA REVISION DOBLE CIEGO.
\autores{}
\adscripcionautor{}

\newcommand{\apellido}{Anonymous et al.}
\newcommand{\titu}{Signal representation and load transfer in bearing fault classification}
\newcommand{\tipo}{Art\'iculo Cient\'if{}ico / Scientif{}ic Paper}
\newcommand{\autor}{}
\newcommand{\pc}{}
\newcommand{\kw}{}
\usepackage{xcolor}
\usepackage{color}
\input{configuracion_paquetes}

\newenvironment{tabla}
  {\par\bigskip\noindent\minipage{\linewidth}}
  {\endminipage\par\bigskip}
\newenvironment{figura}
  {\par\bigskip\noindent\minipage{\linewidth}}
  {\endminipage\par\bigskip}
%========================================================
\begin{document}
\setcounter{page}{1}
% Etiquetas en INGLES, no en espanol como trae la plantilla.
%
% `plantilla_ingenius.tex` (lineas 143-144) y `ingenius.cls` (linea 33, opcion
% `tablename=Tabla` del paquete caption) fijan las etiquetas en espanol, y de
% ahi salian copiadas. Pero la plantilla esta escrita en espanol porque el
% ejemplo lo esta, no porque la revista lo imponga: comprobado el 03/09 sobre
% un articulo en ingles publicado por la propia Ingenius, Contreras Urgiles et
% al. (n.o 21, 2019, DOI 10.17163/ings.n21.2019.03), que es ademas la cita [9]
% de este manuscrito. Sus pies dicen "Figure 1." y "Table 1.", y su
% bibliografia se titula "References".
%
% Antes quedaba una mezcla dentro de la misma frase: el cuerpo decia
% "Table 2 summarises the four" y el pie de esa tabla decia "Tabla 2".
\renewcommand{\refname}{References}
\renewcommand{\figurename}{Figure}
\renewcommand{\tablename}{Table}
\frenchspacing

\begin{tikzpicture}[overlay]
  \node at (138.5mm,12.2mm){\includegraphics[scale=0.35]{figures/logoingenius.pdf}};
  \node at (138mm,2mm){\textsf{\scriptsize pISSN: 1390-650X / eISSN: 1390-860X}};
  \node at (79mm,10mm){\includegraphics[scale=0.6]{figures/creativec.pdf}};
\end{tikzpicture}

\input{silabeo}
\membrete
\input{configuracion_portada}

\begin{textblock}{12.4}(0,10)
\adscripcion
\end{textblock}
%========================================================
\begin{table}[htb]
\begin{tabular}{p{8cm}p{8cm}}
{\Large \bfseries Abstract} &{\Large \bfseries Resumen} \\ \\
Comparisons of signal representations for bearing fault diagnosis are usually
reported under a single operating condition, where accuracy saturates. This
work measures how four representations behave when the test load differs from
the training load. Time-domain statistics, FFT magnitude, db4 wavelet band
energies and STFT spectrogram were combined with three vector classifiers over
the sixteen (training load, test load) pairs of the CWRU drive-end dataset at
12 kHz, with the split made by source record and purged contiguous blocks.
Within a single load, where the split shares records, all four exceeded 0.95 macro-F1 on average. Across loads the mean drop ranged from 0.074 for the spectrogram to
0.153 for time-domain statistics, and the latter depended on transfers involving 0 HP. With default hyperparameters the
omnibus Friedman test did not reject, which does not show that the
representations are equivalent, while the three per-classifier tests each
rejected and disagreed on the ranking. The representation by classifier
interaction persisted when each load was left out and under nested tuning on
training data, which also made the omnibus test reject, so rankings depend on
the configuration. With four loads and one bearing per fault condition these
inferences are indicative and specific to this dataset. Per-window RMS normalisation gave no clear improvement in the four combinations
tested and degraded transfer with Random Forest.
&
Las comparaciones entre representaciones de señal para el diagnóstico de fallas
en rodamientos suelen reportarse bajo una sola condición de operación, donde la
exactitud satura. Este trabajo mide cómo se comportan cuatro
representaciones cuando la carga de prueba difiere de la de entrenamiento. Se
combinaron estadísticos temporales, magnitud FFT, energías de banda wavelet db4
y espectrograma STFT con tres clasificadores de vector sobre los dieciséis pares de cargas del conjunto CWRU a 12 kHz, con
división por registro y por bloques contiguos purgados. Dentro de una
misma carga, donde la partición comparte registros, las cuatro superaron en promedio 0,95 de
F1-macro. Entre cargas la caída media osciló entre 0,074 para el
espectrograma y 0,153 para los estadísticos, y esta
última dependió de las transferencias con 0 HP. Con hiperparámetros por
defecto la prueba de Friedman ómnibus no rechazó, lo que no demuestra su
equivalencia, y las tres pruebas por clasificador
rechazaron y discreparon en el orden. La interacción representación-clasificador
persistió al excluir cada carga y con un ajuste anidado sobre datos de
entrenamiento, que hizo rechazar la prueba ómnibus, por lo que el orden
depende de la configuración. Con cuatro cargas y un rodamiento por condición de
falla, estas inferencias son indicativas y propias de este conjunto. La normalización RMS no mejoró claramente ninguna de las cuatro
combinaciones probadas y degradó la transferencia con Random Forest.
\\[-0.3cm]
\keywords{bearing fault diagnosis, data leakage, domain shift, Friedman test,
signal representation, vibration analysis}
&
\palabrasclave{análisis de vibraciones, desplazamiento de dominio, diagnóstico
de fallas en rodamientos, fuga de datos, prueba de Friedman, representación de
señal}
\\
\end{tabular}
\end{table}

\newpage

\begin{multicols}{2}
\raggedcolumns
% Sin esto, multicols intenta igualar la altura de las dos columnas
% estirando el espacio elastico de las cabeceras de seccion (\subsection lleva
% "plus" en su salto), lo que produce huecos grandes justo despues de 2.2, 2.5,
% 3.3 y otras subsecciones. \raggedcolumns deja el pie de columna irregular en
% vez de forzar ese estiramiento. Detectado visualmente en el PDF: hay que
% revisar siempre el render, "cero Overfull" no basta.
%========================================================
\section{Introduction}

Rolling element bearings fail in ways that leave a signature in the vibration
measured on the housing: a localised defect produces an impact each time a
rolling element passes over it, exciting the structural resonance at a rate set
by the geometry and the shaft speed, which is the basis of envelope analysis
\cite{randall2011}.

Data-driven diagnosis converts the signal into a representation and trains a
classifier on it, and the representation is normally chosen by comparing
candidates on the same data. The Case Western Reserve University (CWRU) Bearing
Data Center dataset \cite{cwru} is the standard reference for those experiments,
with a benchmark study devoted to it \cite{smith2015} and later work refining
that practice \cite{hendriks2022}.

Accuracy within a single operating condition saturates near unity, and the way
the data are partitioned decides the result: bearing vibration is usually cut
into overlapping windows, and if those are shuffled before splitting, windows
sharing samples fall on both sides and the classifier recognises the recording
rather than the fault. Data leakage was found in 294 papers across
seventeen scientific fields, in some cases producing conclusions far more
optimistic than the data supported \cite{kapoor2023}.

The shift between operating conditions is itself an active subject. The CWRU
dataset is a common validation reference \cite{neupane2020}, generalisation
remains an open problem across nine datasets \cite{zhao2020}, and the usual
response is a domain-adaptation model aligning a labelled source condition and
an unlabelled target condition, for instance with a multi-kernel maximum mean
discrepancy \cite{li2019}. That line fixes a representation and modifies the
model, leaving open what the representation itself contributes to the transfer,
which is the question addressed here.

Fault diagnosis in rotating machinery has been addressed repeatedly in this
journal from measurement rather than model design: time-domain descriptors fed a feed-forward network to diagnose spark-ignition
engines \cite{contreras2019}; temperature and injection pressure served as measured indicators analysed by multifactorial ANOVA \cite{llanes2019}; two measurement methods for stator lamination faults
were compared without a learned model \cite{salazar2012}; and a four-level
Daubechies db4 decomposition was used to estimate fault distance
\cite{gomez2023}, the family used below with five levels.

This study measures how four representations behave when the test load differs
from the training load. The contributions of this work are: (1) a
record-level evaluation protocol for the CWRU dataset in which the split is
performed by source record and, within a record, by purged contiguous blocks,
verified by a leakage check with a negative control; (2) a 4$\times$4 matrix of
macro-F1 by (training load, test load) pair for four representations and three
classifiers, reported per fold; (3) evidence, specific to this dataset and to
the configurations evaluated, that the ranking of representations depends on the
classifier, with a representation by classifier interaction that persists when
any load is left out and under nested tuning; and (4) an ablation, specified
before the experiment was run, in which per-window RMS division did not clearly improve
cross-load transfer with the two classifiers and representations tested.

\section{Materials and Methods}

\subsection{Dataset and label verification}

Forty drive-end accelerometer records of the CWRU Bearing Data Center
\cite{cwru}, sampled at 12 kHz, cover ten class-diameter combinations under four
loads (0, 1, 2 and 3 HP, nominally 1797, 1772, 1750 and 1730 rpm). The classes
are healthy bearing (Normal), inner race fault (IR), ball fault (B) and outer
race fault (OR), the last three at diameters of 0.007, 0.014 and 0.021 inches.

The file naming convention is a known source of error: \texttt{99.mat} (Normal,
2 HP) also holds a duplicate of the 1 HP signal, so the channel is matched
exactly against the file number (\texttt{X\{NNN\}\_DE\_time}).

Labels were verified against the physics of the SKF 6205 drive-end bearing
rather than third-party tables, by envelope analysis of the 2--5.5 kHz
resonance at its characteristic orders; the rolling-element value of 4.7135
tabulated by CWRU corresponds to $2\times$BSF, not to BSF. The inner race labels
were confirmed in 12 of 12 records, the outer race labels in 9 of 12 and the
ball labels, by the $2\times$BSF sideband criterion, in 11 of 12, with weak
signatures in the outer race records at 0.014 in and the ball records at
0.021 in; the details are in the supplementary material. All 40 records were
kept, since discarding those with a weak signature would bias the result
upwards, and fault diameter is a covariate for subgroup reporting.

\subsection{Windowing and the partition rule}

Each record was segmented into windows of 2048 samples with 50 \% overlap,
yielding 5887 windows. The split is performed by source record, never by window,
because overlapping windows of a record share samples.

Outside the diagonal this holds by construction, since training on load $i$ and
testing on load $j$ uses different files; on the diagonal it cannot, because the
dataset provides a single healthy record per load. One rule was therefore
adopted for all sixteen cells: the windows of each record are cut into five
contiguous blocks in time and the first window of each internal block is
discarded. With a stride of 1024 samples and a window of 2048, windows $m$ and
$m+2$ share no samples, so one discarded window per boundary suffices, removing
160 windows (2.7 \%) and leaving 5727. For cell $(i,j)$ and fold $k$, training
uses the blocks other than $k$ from every record of load $i$, and testing block
$k$ from every record of load $j$.

On the diagonal, overlap is removed but
training and test windows still come from the same record, so the diagonal is a
reference and is excluded from every test.

All 80 cell-fold combinations were checked for training windows sharing samples
with a test window of the same record, and none was found. The same check
without the purge reports 320 violations, which is the 160 overlapping pairs
created by four block boundaries in each of the forty records, each counted in
the two folds it affects. Table \ref{tab:ventanas} gives the usable window
counts; the healthy records last 20 to 40 s against about 10 s for the fault
records.

\begin{table}[H]
\centering
\captionof{table}{Usable windows of 2048 samples with 50 \% overlap, by class
and load, after discarding the boundary windows between blocks.}
\label{tab:ventanas}
\begin{tabular}{lrrrr}
\toprule
Class & 0 HP & 1 HP & 2 HP & 3 HP \\
\midrule
Normal & 233 & 467 & 468 & 469 \\
IR     & 340 & 340 & 340 & 342 \\
B      & 341 & 340 & 341 & 341 \\
OR     & 341 & 342 & 340 & 342 \\
\bottomrule
\end{tabular}
\end{table}

\subsection{Signal representations}

Four representations were computed from the same windows. Figure
\ref{fig:senales} shows one raw window per class and Figure \ref{fig:repr} shows
the same inner-race window under the four representations.

\textbf{Time-domain statistics.} Twelve classical vibration descriptors, from mean, RMS and kurtosis to the crest, shape, impulse and clearance factors.

\textbf{FFT magnitude.} The real FFT magnitude up to Nyquist, 1025 bins of
5.86 Hz.

\textbf{Wavelet db4.} Five-level discrete decomposition with the Daubechies db4
wavelet, the family used in \cite{gomez2023}, as the relative energy of each
band (six values: the approximation A5 and the details D5 to D1).

\textbf{STFT spectrogram.} The spectrogram in dB, normalised per image to
$[0,1]$, with $\mathit{nperseg} = 64$: 187.5 Hz per band, 63 frames and a hop of
2.67 ms. A fault excites the 2--5.5 kHz structural resonance, and its
characteristic frequency appears as the repetition rate of the impacts, every
9.3 ms (BPFO) and 6.2 ms (BPFI) at 1797 rpm, so the hop places 2.3 to 3.5
frames within each impact period. The comparison uses a $16\times64$ grid of
1024 values, which keeps the hop and halves the frequency axis to 375 Hz; a
square $32\times32$ grid would leave fewer than two columns per BPFI period. A
$64\times64$ image of the same transform feeds the convolutional network.

Table \ref{tab:repr} summarises the four. The scaler is fitted inside the
pipeline, on the training fold only.

\begin{table}[H]
\centering
\captionof{table}{Dimension, effective resolution and dependence on absolute
amplitude of each representation.}
\label{tab:repr}
\resizebox{\columnwidth}{!}{
\begin{tabular}{lrlcc}
\toprule
Representation & Dim. & Effective resolution & Amp.-inv. & Enters comparison \\
\midrule
Time-domain statistics       & 12   & ---                  & no  & yes \\
FFT magnitude                & 1025 & 5.86 Hz              & no  & yes \\
Wavelet db4 (rel. energy)    & 6    & dyadic band          & yes & yes \\
STFT spectrogram ($16\times64$) & 1024 & 375.0 Hz / 2.67 ms & yes & yes \\
STFT spectrogram ($64\times64$) & 4096 & 187.5 Hz / 2.67 ms & yes & no  \\
\bottomrule
\end{tabular}}
\end{table}

Two of the four discard amplitude by construction, the wavelet through relative
band energy and the spectrogram through per-image normalisation, while the other
two keep the absolute scale, which varies between loads. An ablation was therefore
fixed before the experiment was run: time-domain statistics and FFT magnitude
recomputed on windows divided by their RMS value, with SVM-RBF and Random Forest
over the twelve off-diagonal pairs.

\subsection{Classifiers}

Three classifiers that accept vectors were used with all four representations:
SVM with RBF kernel, Random Forest with 300 trees and a two-layer perceptron of
128 and 64 units trained for at most 500 iterations, all from scikit-learn
\cite{pedregosa2011}. Class weighting was balanced for the SVM and the forest;
the perceptron has no class-weight option and was trained unweighted. A standard scaler was
fitted on the training fold for the SVM and the perceptron, and omitted for the
forest. Otherwise, hyperparameters were left at the library defaults in the protocol, so
that all four representations meet each classifier under identical settings.
For the SVM this fixes $C = 1$ and $\gamma = 1/(d\,\mathrm{Var}(X))$, which runs
from $1/1025$ for FFT magnitude to $1/6$ for the wavelet energies.

As a sensitivity analysis added in revision, each classifier was also tuned by
nested search on training data only: for each training load and fold, each of
the four training blocks was held out in turn and the configuration with the
best mean inner macro-F1 was refitted on all four, without the test load or
block. The 12 SVM, 6 forest and 6 perceptron configurations, the tie rule and
the choice in each fold are in the supplementary material. Since the inner
validation stays within one load, where macro-F1 saturates, the search
optimises in-domain performance, not transfer.

A two-dimensional convolutional network provides a reference row: three blocks
of 32, 64 and 128 channels with batch normalisation and max pooling, then global
average pooling and a linear layer, in PyTorch \cite{paszke2019}, trained on
$64\times64$ spectrograms for a fixed 30 epochs with Adam at $10^{-3}$ and a
batch of 64. The epoch count was not tuned, and the network is a descriptive reference only: it has no
counterpart for the other three representations, and an unbalanced design
invalidates both Friedman and the ANOVA.

\subsection{Experimental design and statistical analysis}

The design covers 4 representations $\times$ 3 vector classifiers $\times$ 16
cells $\times$ 5 folds, that is 960 fits, plus 80 for the convolutional
reference and 240 for the ablation. The seed was fixed at 42 and results stored
per fold. 

Analyses differ in status. The protocol, fixed on 20 August 2026 before the data
were split, specified a Friedman test over the four representations with the
off-diagonal pairs as blocks, a Nemenyi post-hoc test and a two-way ANOVA of
representation by load distance. On 22--23 August, before the experiment, the
primary test was set to macro-F1 averaged over classifiers and folds ($N = 12$,
$k = 4$) \cite{demsar2006}, one Friedman test per classifier was added at a
Bonferroni-corrected $\alpha = 0.0167$, and the amplitude ablation was fixed.
Only these are confirmatory; every later analysis is exploratory, with
uncorrected $p$-values.

Averaging the classifiers keeps $N$ at 12, since counting each (pair,
classifier) as a block would rest on their shared partition. The pairs are not
independent either: each load appears in six, and $(i, j)$ and $(j, i)$ reuse
the same records. The effective $N$ is below 12, so the Friedman test is
anticonservative and the Nemenyi critical distance too short, which weakens the
rejecting per-classifier tests more than the non-rejecting primary one. Two analyses added in revision address this dependence. A linear mixed model on the cell means
models it, with training load and test load as crossed random effects and their
interactions with representation and with representation by classifier, tested
by likelihood ratio with a 999-replicate parametric bootstrap. A
leave-one-load-out jackknife of the Friedman tests, the interaction and the drops
measures how much each result depends on a single load. With four loads no inferential procedure is fully
reliable, so all $p$-values are indicative.

The two-way ANOVA of the protocol treats folds as replicates although they share
four fifths of their training windows, so its degrees of freedom, $p$-values and
effect sizes are not interpreted. It was refitted on the 48 cell means, and a
three-way model on the 144 (representation, classifier, pair) cell means adds
the classifier, with a variant blocking on the pair. These cell means are not
independent either, so the residual degrees of freedom are overstated and the
$p$-values of these models are anticonservative. The ablation is reported, as in
the protocol, as a paired difference over the 60 fold-level differences with a
bootstrap interval and Cliff's delta; since folds share training windows, the
intervals were also computed over the twelve pairs (supplementary material).

\begin{figure}[H]
\begin{center}
\includegraphics[width=8.17cm]{figures/fig1_p7.pdf}
\end{center}
\caption{Raw drive-end acceleration at 0 HP, one 2048-sample window per class:
(a) Normal, (b) inner race, (c) ball, (d) outer race. Each panel has its own
$y$-axis scale.}
\label{fig:senales}
\end{figure}

\begin{figure}[H]
\begin{center}
\includegraphics[width=8.17cm]{figures/fig2_p7.pdf}
\end{center}
\caption{Same inner-race window at 0 HP under the four representations:
(a) time-domain statistics, (b) FFT magnitude, (c) db4 wavelet relative energy
per band, from A5 to D1, (d) STFT spectrogram on the $64\times64$ CNN grid, not
the $16\times64$ grid that enters the comparison.}
\label{fig:repr}
\end{figure}

\section{Results and Discussion}

\subsection{In-domain performance saturates}

Within a single load all four representations exceed 0.95 macro-F1, averaged
over classifiers: 0.9997 for FFT magnitude, 0.9925 for the spectrogram, 0.9882
for time-domain statistics and 0.9602 for the wavelet, with the breakdown by
classifier in Table \ref{tab:f1}. Under four hundredths separate best from worst
and ten of the twelve representation-classifier combinations sit at or above
0.98, so on this dataset a comparison restricted to the training condition
separates little.

\begin{table}[H]
\centering
\captionof{table}{Macro-F1 (mean $\pm$ standard deviation over folds and cells)
in-domain and out-of-domain, by representation and classifier, with default
hyperparameters. Bold marks the best value per column among the twelve compared
rows; the in-domain column has two tied at 1.000. The convolutional network,
below the line, is a descriptive reference and enters no comparison.}
\label{tab:f1}
\resizebox{\columnwidth}{!}{
\begin{tabular}{llcc}
\toprule
Representation & Classifier & F1 in-domain & F1 out-of-domain \\
\midrule
Statistical & MLP           & 0.988 $\pm$ 0.015 & 0.799 $\pm$ 0.229 \\
Statistical & Random Forest & 0.992 $\pm$ 0.008 & 0.869 $\pm$ 0.103 \\
Statistical & SVM-RBF       & 0.985 $\pm$ 0.020 & 0.836 $\pm$ 0.174 \\
FFT & MLP                   & \textbf{1.000 $\pm$ 0.000} & \textbf{0.971 $\pm$ 0.043} \\
FFT & Random Forest         & \textbf{1.000 $\pm$ 0.000} & 0.956 $\pm$ 0.054 \\
FFT & SVM-RBF               & 0.999 $\pm$ 0.003 & 0.774 $\pm$ 0.061 \\
Wavelet db4 & MLP           & 0.969 $\pm$ 0.017 & 0.830 $\pm$ 0.087 \\
Wavelet db4 & Random Forest & 0.982 $\pm$ 0.011 & 0.892 $\pm$ 0.045 \\
Wavelet db4 & SVM-RBF       & 0.929 $\pm$ 0.018 & 0.850 $\pm$ 0.074 \\
STFT spectrogram & MLP           & 0.998 $\pm$ 0.003 & 0.916 $\pm$ 0.056 \\
STFT spectrogram & Random Forest & 0.980 $\pm$ 0.011 & 0.898 $\pm$ 0.069 \\
STFT spectrogram & SVM-RBF       & 0.999 $\pm$ 0.001 & 0.942 $\pm$ 0.054 \\
\midrule
\multicolumn{4}{l}{\textit{Descriptive reference, not compared}} \\
CNN 2D spectrogram & CNN2D  & 0.971 $\pm$ 0.116 & 0.957 $\pm$ 0.119 \\
\bottomrule
\end{tabular}}
\end{table}

\subsection{Cross-load degradation}

Figure \ref{fig:mapa} shows macro-F1 for every (training load, test load) pair,
one panel per representation, diagonal boxed. Averaged over the twelve
off-diagonal cells, the drop from the diagonal is 0.0740 for the spectrogram,
0.0995 for FFT magnitude, 0.1029 for the wavelet and 0.1534 for time-domain
statistics. These are numerical differences. Resampling the twelve pairs gives intervals of 0.043--0.105, 0.080--0.122,
0.073--0.133 and 0.074--0.240, approximate because the pairs share loads; the
primary test reported below does not separate the four representations, and
leaving one load out moves the drops within 0.052--0.096, 0.074--0.113,
0.067--0.132 and 0.024--0.207, reversing the position of time-domain statistics.
No ordering is statistically supported, and the smaller drop
of the spectrogram is not evidence that it is intrinsically more robust to load
changes. Descriptively, the large drop of time-domain statistics comes from
transfers involving 0 HP: without that load it is 0.024, the smallest of the
four.



The matrix is also asymmetric. With time-domain statistics, training at 0 HP and
testing at 1 HP gives 0.795 against 0.678 for the reverse; at maximum separation
the two directions give 0.771 and 0.539, the worst cell of the matrix. With
statistics and the wavelet, testing at 0 HP is harder in every pair, and the
healthy class has 233 windows there against about 468 elsewhere, while FFT
magnitude shows the opposite asymmetry. The healthy class does not account for all of it. Averaged over pairs, macro-F1 transferring towards the lower load minus
towards the higher is $-0.082$ for time-domain statistics and $-0.063$ with the
healthy class excluded, and the other three representations also keep their sign
without it.

The convolutional network is a descriptive reference only (Table \ref{tab:f1}):
it uses a different input and architecture and enters no comparison. Trained at
3 HP it failed in one of the five folds, whose four cells fall between 0.396 and
0.489 and produce its standard deviations of 0.116 and 0.119; it is reported as
it ran, since the epoch budget was fixed in advance.

% Fig. 3 es mas ancha que una columna (16 cm) y no cabe a media columna, asi
% que no puede ir con [H] dentro del flujo de multicols como las demas figuras
% y tablas. Como figure*[t] normal, LaTeX la trata como flotante en cola: espera
% el siguiente hueco de pagina completa, que caia DESPUES de donde el texto ya
% habia colocado la Fig. 4 (columna simple, [H]), invirtiendo el orden de
% lectura. Probado \clearpage antes de esta figura: empeora el problema (deja
% una pagina casi vacia y aun asi no garantiza el orden, porque clearpage solo
% termina la pagina, no fuerza que ESTE flotante en particular ocupe la
% siguiente). La solucion correcta es sacarla de multicols: se cierra el
% entorno de dos columnas, se inserta la figura a ancho completo exactamente
% aqui (sin cola de flotantes, sin ambiguedad de orden) y se reabre multicols
% despues. Es la tecnica estandar para figuras anchas dentro de multicols.
\end{multicols}
\begin{figure}[H]
\begin{center}
\includegraphics[width=16cm]{figures/fig3_p7.pdf}
\end{center}
\caption{Macro-F1 for every (training load, test load) pair, averaged over the
three vector classifiers and 5 folds. Rows are the training load and columns the
test load. (a) Statistical, (b) FFT, (c) Wavelet db4,
(d) STFT spectrogram. The diagonal is a reference cell, boxed for orientation,
and takes part in no test.}
\label{fig:mapa}
\end{figure}
\begin{multicols}{2}
\raggedcolumns

\subsection{No consistent ranking across classifiers}

With the default settings of the protocol, the primary Friedman test over the
twelve off-diagonal pairs does not reject ($\chi^2 = 5.20$, $p = 0.158$; mean
ranks 1.833 for the spectrogram, 2.500 for FFT magnitude, 2.667 for time-domain
statistics and 3.000 for the wavelet, against a Nemenyi critical distance of
1.354), and no pair is separated in Figure \ref{fig:nemenyi}, the smallest
post-hoc $p$-value being 0.119.

The three per-classifier tests each reject at the Bonferroni-corrected
$\alpha = 0.0167$ and disagree on the ordering. The spectrogram ranks first with
SVM-RBF ($\chi^2 = 17.10$, $p = 0.00067$), where FFT magnitude ranks last at
3.500, while FFT magnitude ranks first with Random Forest ($\chi^2 = 13.70$,
$p = 0.0033$, rank 1.333) and the perceptron ($\chi^2 = 21.00$, $p = 0.00011$,
rank 1.083). Averaging them can offset that disagreement, which is consistent with the omnibus
test not rejecting while its components do.

The fold-level ANOVA of the protocol is not interpreted. On the 48 cell means
only load distance remains ($F = 5.85$, $p = 0.006$); representation gives
$F = 2.62$, $p = 0.065$, and its interaction with load distance $F = 0.85$,
$p = 0.543$. Adding the classifier locates the effect elsewhere (Table
\ref{tab:anova}): in these configurations the representation by classifier
interaction is the largest term ($F = 4.91$, $p < 0.001$, $\eta^2_p = 0.214$),
while representation by load distance is not significant ($F = 1.76$,
$p = 0.113$). With the pair as a block the interaction stays at $F = 5.78$
($p < 0.001$) and the block has the largest partial effect size ($\eta^2_p = 0.307$) in this
exploratory model.

The analyses added for the dependence between pairs qualify these results
without overturning them. In the jackknife the interaction remains in all four
reduced sets ($F$ from 2.72 to 11.66, $p \le 0.022$) and each per-classifier
leader is unchanged (once tied), although with $N = 6$ eight of the twelve tests
fall short of the corrected threshold (Kendall's $W$ from 0.30 to 0.94). The
primary ordering is less stable: without 0 HP time-domain statistics rank first
and the test rejects ($p = 0.006$), and without any other load they rank last.
The mixed model could be fitted only without the classifier by load terms, and
a quarter or more of its bootstrap refits failed, so its likelihood ratios, 2.29
for representation and 21.74 for the interaction (indicative $p = 0.61$ and
$0.005$), are not relied on.

Nested tuning changes the rankings but not the disagreement. The per-classifier
tests still reject with the same leaders, the blocked interaction stays at
$F = 4.66$ ($p < 0.001$), and the primary test now rejects ($\chi^2 = 9.30$,
$p = 0.026$; mean ranks 1.92 for FFT magnitude, 2.17 for the spectrogram, 2.50
for the statistics and 3.42 for the wavelet). The change comes mostly from the
SVM: tuning raised its cross-load macro-F1 with FFT magnitude from 0.774 to
0.833, out of last place, and with the wavelet raised in-domain macro-F1 from
0.929 to 0.975 while cutting cross-load macro-F1 from 0.850 to 0.680, consistent with an inner validation that rewards fit to the training load. The forest kept its defaults
in 69 of 80 models. The rankings are therefore specific to the configurations
evaluated, not the best each representation can reach, and whether the primary test rejects depends on them.

\begin{table}[H]
\centering
\captionof{table}{Exploratory three-way ANOVA on macro-F1 over the 144
off-diagonal cell means, with default hyperparameters; $p$-values are not
adjusted for the dependence between cells.}
\label{tab:anova}
\resizebox{\columnwidth}{!}{
\begin{tabular}{lrrrrr}
\toprule
Source & SS & df & $F$ & $p$ & $\eta^2_p$ \\
\midrule
Representation (R)      & 0.1595 & 3   & 5.47  & 0.002    & 0.132 \\
Classifier (C)          & 0.0687 & 2   & 3.53  & 0.033    & 0.061 \\
Load distance (D)       & 0.2370 & 2   & 12.18 & $<0.001$ & 0.184 \\
R $\times$ C            & 0.2865 & 6   & 4.91  & $<0.001$ & 0.214 \\
R $\times$ D            & 0.1030 & 6   & 1.76  & 0.113    & 0.089 \\
C $\times$ D            & 0.0178 & 4   & 0.46  & 0.766    & 0.017 \\
R $\times$ C $\times$ D & 0.0336 & 12  & 0.29  & 0.990    & 0.031 \\
Residual                & 1.0504 & 108 &       &          &       \\
\bottomrule
\end{tabular}}
\end{table}

\begin{figure}[H]
\begin{center}
\includegraphics[width=8.17cm]{figures/fig4_p7.pdf}
\end{center}
\caption{Nemenyi critical-difference diagram for the primary Friedman test
($N = 12$ off-diagonal pairs, $k = 4$ representations, $\chi^2 = 5.20$,
$p = 0.158$). Representations joined by a bar are not distinguishable at
$\alpha = 0.05$; CD is the Nemenyi critical distance.}
\label{fig:nemenyi}
\end{figure}

\subsection{Performance by fault diameter}

Over the off-diagonal cells recall is highest for the healthy class (0.994) and
lowest for inner race faults at 0.014 in (0.587). The next four subgroups are
ball at 0.021 in (0.786), inner race at 0.021 in (0.823), inner race at 0.007 in
(0.832) and outer race at 0.014 in (0.849); the remaining four sit above 0.89.
For inner race faults at 0.014 in recall is 0.836 with FFT magnitude against
0.485 with the spectrogram, 0.494 with time-domain statistics and 0.533 with the
wavelet. These figures describe where classification fails across loads, not why: label
verification found weak signatures in two of the weaker subgroups but not in the
weakest, and nothing here identifies a physical cause. The subgroup analysis
is descriptive and secondary to the question of the study.

\subsection{Amplitude normalisation in the tested configurations}

The hypothesis behind the ablation is that representations retaining absolute
amplitude degrade because amplitude scales with load. Dividing each window by
its RMS value before extracting time-domain statistics and FFT magnitude tests
that procedure with SVM-RBF and Random Forest. No combination improved by an
interval that excludes zero, and two degraded. With Random Forest, time-domain
statistics lost 0.149 macro-F1 (95 \% CI $-0.197$ to $-0.103$, Cliff's delta
$-0.472$) and FFT magnitude lost 0.133 (95 \% CI $-0.151$ to $-0.116$, delta
$-0.750$). With SVM-RBF both intervals cross zero: $-0.018$ for statistics
(95 \% CI $-0.047$ to $+0.011$, delta $-0.205$) and $+0.033$ for FFT (95 \% CI
$-0.002$ to $+0.067$, delta 0.330). The result separates the two classifiers
rather than ruling out any benefit. It does not show that amplitude
normalisation is ineffective in general, only that this procedure gave no
consistent improvement here; on this dataset, removing amplitude this way is
inadvisable with the forest and of uncertain effect with the SVM.

\subsection{Limitations}

Each combination of fault type and diameter in the CWRU dataset is a single
physical bearing measured under four loads, so crossing loads keeps the bearing
fixed and a model may rely on the signature of that specimen rather than on the
fault type. The results measure transfer between loads of the same specimens,
not generalisation to other bearings or machines. For that purpose they are
likely optimistic, since specimen-specific cues may survive a change of load, and
they should not be carried to industrial settings without external validation.

The diagonal cells cannot use a record-level split, since the dataset provides
one healthy record per load, so the in-domain figures of Table \ref{tab:f1} are
a reference level, not an estimate of generalisation. The same record leaves
the healthy class imbalanced differently across loads, 233 windows at 0 HP
against about 468 elsewhere; macro-F1 and class weighting mitigate this without
removing the influence of the dataset structure, and the perceptron accepts no
class weights. The drop of time-domain statistics is
concentrated in transfers involving 0 HP, although the direction asymmetry is not confined to the healthy class.

The envelope spectrum, used here to verify the labels, is not among the four representations, which limits
what the comparison covers.

Tuning used only the training load, so the tuned results are not the best
transfer each representation can reach, and the grids are small. The four
representations also differ in dimension from 6 to 1025, which the default
$\gamma$ of the SVM ties to the kernel width; dimension was not equalised
because it is part of a representation.

With four loads every inference rests on few independent units: the mixed model
could not be fitted in full, its bootstrap was unstable, and the jackknife
leaves six pairs. The $p$-value of 0.158 is a failure to separate the four
representations, not evidence that they are equivalent.

Four limitations therefore remain unresolved: generalisation beyond one
specimen per fault condition, which bounds every result; inference with four
loads, which makes all $p$-values indicative; the dependence of the rankings on
the configuration, which restricts them to the settings evaluated; and the cause
of the weak subgroups, which these data do not identify.

\section{Conclusions}

Within a single load the four representations average above 0.95 macro-F1 and
separate by less than four hundredths, so on this dataset a comparison under the
training condition alone does not identify a preferable one. Across loads the
mean drops range from 0.0740 for the STFT spectrogram to 0.1534 for time-domain
statistics, but the primary test does not separate the representations and the largest depends on 0 HP, so
no ordering is established.

No representation was superior across the classifiers and configurations
evaluated. With default hyperparameters the omnibus Friedman test does not
reject, which is not evidence of equivalence, while the per-classifier tests
reject and place different representations first, FFT magnitude ranking last
with SVM-RBF and first with the other two. The representation by classifier
interaction persists when any load is left out and under nested tuning, which
however changes the rankings and makes the omnibus test reject, so it is
exploratory evidence specific to these configurations: on this dataset a
representation is best assessed with the classifier it will be used with.

Per-window RMS normalisation gave no clear improvement in cross-load transfer in
the four tested combinations and degraded it with Random Forest, by 0.149 for time-domain
statistics and 0.133 for FFT magnitude.

All of this concerns transfer between loads of the same bearings. External validation on other bearings and machines is needed before industrial use.

%========================================================
\section*{Data availability}

% VERSION PARA REVISION DOBLE CIEGO: el identificador del repositorio se
% retiene. Texto para la version aceptada: ver ENVIO.md del proyecto.
The CWRU records are public \cite{cwru}. The repository, withheld for
double-blind review and provided upon acceptance, holds the per-fold results and
the code of every step, from download, label verification, partition and
leakage check to feature extraction, classifier configuration, metrics,
statistics, nested tuning and figures. One command reruns the pipeline from the
raw records, whose download by the first script is the only external step, and
regenerates every table and figure. On CPU the results are identical with the
fixed seed, although the perceptron on low-dimensional inputs also depends on the number of BLAS threads (the pipeline uses the default), and the
convolutional network reproduces on the original GPU. A full rerun from a fresh copy of the repository
regenerated every table identically, apart from timing columns.
%========================================================

\bibliography{referencias.bib}
\bibliographystyle{IEEEtran}
\end{multicols}
\end{document}
